\section{Chi tiết Pipeline (Jenkinsfile)}

Pipeline được viết bằng \textbf{Declarative Pipeline} trong file \texttt{Jenkinsfile} tại root của repository. Tổng cộng 298 dòng, 8 stages chính.

\subsection{Cấu hình chung}
\begin{lstlisting}[language=java, caption={Cấu hình pipeline}]
pipeline {
    agent any
    tools {
        maven 'Maven'
        jdk   'JDK-25'
    }
    options {
        timestamps()
        buildDiscarder(logRotator(numToKeepStr: '10'))
        timeout(time: 60, unit: 'MINUTES')
        disableConcurrentBuilds()
    }
    environment {
        COVERAGE_THRESHOLD = '70'
        SONAR_HOST_URL     = 'http://3.27.92.213:9000'
        MAVEN_OPTS         = '-Xmx512m -XX:MaxMetaspaceSize=256m'
    }
}
\end{lstlisting}

Các cấu hình quan trọng:
\begin{itemize}[noitemsep]
    \item \texttt{COVERAGE\_THRESHOLD = '70'}: Ngưỡng coverage tối thiểu 70\%
    \item \texttt{MAVEN\_OPTS}: Giới hạn bộ nhớ JVM để phù hợp với EC2
    \item \texttt{disableConcurrentBuilds()}: Tránh xung đột khi nhiều builds chạy cùng lúc
\end{itemize}

\subsection{Stage 1: Checkout}
\begin{lstlisting}[language=java, caption={Stage Checkout}]
stage('Checkout') {
    steps {
        checkout scm
        sh 'git fetch origin main:refs/remotes/origin/main || true'
        sh 'chmod +x ci/detect-changed-modules.sh ci/check-coverage.sh'
    }
}
\end{lstlisting}

Stage này thực hiện:
\begin{itemize}[noitemsep]
    \item Clone source code từ GitHub
    \item Fetch \texttt{origin/main} để có baseline cho việc so sánh thay đổi (fix shallow clone của Jenkins)
    \item Cấp quyền thực thi cho các CI scripts
\end{itemize}

\subsection{Stage 2: Detect Changed Modules}
\begin{lstlisting}[language=java, caption={Stage Detect Changed Modules}]
stage('Detect Changed Modules') {
    steps {
        script {
            env.CHANGED_MODULES = sh(
                script: 'ci/detect-changed-modules.sh',
                returnStdout: true
            ).trim()
        }
    }
}
\end{lstlisting}

Stage này gọi script \texttt{ci/detect-changed-modules.sh} để xác định modules nào thay đổi so với \texttt{origin/main}. Chi tiết về script được mô tả ở Mục \ref{sec:monorepo}.

\subsection{Stage 3: Gitleaks -- Secret Scan}
Quét toàn bộ source code tìm secrets/API keys bị rò rỉ. Chi tiết ở Mục \ref{sec:gitleaks}.

\subsection{Stage 4: Test (JUnit + JaCoCo)}
\begin{lstlisting}[language=java, caption={Stage Test}]
stage('Test') {
    steps {
        script {
            def modules = env.CHANGED_MODULES.split(',')
            for (module in modules) {
                sh "mvn -B -ntp -pl ${module} -am test jacoco:report -DskipITs"
            }
        }
    }
    post {
        always {
            junit allowEmptyResults: true,
                  testResults: '**/target/surefire-reports/*.xml'
            archiveArtifacts allowEmptyArchive: true,
                             artifacts: '**/target/site/jacoco/jacoco.xml'
        }
    }
}
\end{lstlisting}

Stage này:
\begin{itemize}[noitemsep]
    \item Chỉ chạy test cho \textbf{modules thay đổi} (\texttt{-pl \$\{module\}})
    \item Tự động build dependencies (\texttt{-am}: also make)
    \item Skip Integration Tests (\texttt{-DskipITs})
    \item Sinh JaCoCo coverage report (\texttt{jacoco:report})
    \item Upload test results (JUnit XML) và coverage report (JaCoCo XML) lên Jenkins
\end{itemize}

\subsection{Stage 5: Coverage Gate}
\begin{lstlisting}[language=java, caption={Stage Coverage Gate}]
stage('Coverage Gate') {
    steps {
        script {
            def modules = env.CHANGED_MODULES.split(',')
            def serviceModules = modules.findAll { it != 'common-library' }
            if (serviceModules.isEmpty()) {
                echo ">>> No service modules to check coverage"
            } else {
                for (module in serviceModules) {
                    sh "ci/check-coverage.sh ${module} ${env.COVERAGE_THRESHOLD}"
                }
            }
        }
    }
}
\end{lstlisting}

Stage này:
\begin{itemize}[noitemsep]
    \item Gọi script \texttt{ci/check-coverage.sh} để kiểm tra coverage $\geq$ 70\%
    \item \textbf{Exclude common-library}: Đây là shared library với coverage 33\% từ upstream, không phải service module
    \item Script parse JaCoCo XML report, tính LINE coverage percentage
    \item Pipeline FAIL nếu bất kỳ service module nào có coverage $<$ 70\%
\end{itemize}

\subsection{Stage 6: Build}
\begin{lstlisting}[language=java, caption={Stage Build}]
stage('Build') {
    steps {
        script {
            def modules = env.CHANGED_MODULES.split(',')
            for (module in modules) {
                sh "mvn -B -ntp -DskipTests -pl ${module} -am package"
            }
        }
    }
}
\end{lstlisting}

Build chỉ các modules thay đổi, skip tests (đã chạy ở stage 4).

\subsection{Stage 7: SonarQube Analysis}
Phân tích chất lượng code. Chi tiết ở Mục \ref{sec:sonarqube}.

\subsection{Stage 8: Snyk Dependency Scan}
Quét lỗ hổng bảo mật trong dependencies. Chi tiết ở Mục \ref{sec:snyk}.

\subsection{Post Actions}
\begin{lstlisting}[language=java, caption={Post-build actions}]
post {
    always { cleanWs() }
    success {
        echo "PIPELINE PASSED SUCCESSFULLY"
    }
    failure {
        echo "PIPELINE FAILED - Check console output"
    }
}
\end{lstlisting}

\begin{itemize}[noitemsep]
    \item \texttt{cleanWs()}: Dọn workspace sau mỗi build để tiết kiệm disk
    \item In thông báo kết quả (pass/fail) kèm branch name và build number
\end{itemize}
